% jsc2019_qual B - Kleene Inversion の手書きメモの再現と振り返り
% ビルド: tools/build_thinking.sh jsc2019_qual/b --preview
\documentclass[a4paper]{article}
\usepackage[margin=1.2cm]{geometry}
\usepackage{handmemo}
\usepackage{hyperref}
\pagestyle{empty}

\begin{document}

% ---------------------------------------------------------------- 1 枚目
\begin{memopage}{20260925\_174814.jpg}
  \hw(1.5,1.8){B --- Kleene\ \ Inversion}

  \hw(1.5,3.3){長さ N}
  \hw(4.4,3.3){$A = [A_0, \cdots, A_{N-1}]$}

  \hwm(1.5,5.1){K \times \underset{A}{\underset{\sim}{N}} : B}
  \redpen(6.8,5.0){K は $10^9$ まで、N は 2000 まで。\\B は作らずに数える}

  \hw(1.5,7.1){$A = 1, 3, 2,$\ \ $K = 2$}
  \hw(1.5,8.7){$B = 1, 3, 2, 1, 3, 2$}

  \hw(1.5,10.5){Bの転倒数\ \ $t$\ \ $\%\ (10^9 + 7)$}

  \hw(1.4,12.8){$t : \mathrm{count}(\ (i, j)\ \ (0 \leq i < j \leq K \times N - 1) \wedge B_i > B_j\ )$}
  \redpen(8.6,10.3){\small 同じコピー内: $K \cdot x$ 個（$x$ = A の中の $i<j,\ A_i>A_j$）\\
    \small コピー間: 位置に関係なく $A_i > A_j$ の組 × $\frac{K(K-1)}{2}$}
  \redarrow(11.6,11.2)(10.6,12.3)

  % サンプル 1
  \hw(0.9,15.5){2\ 2}
  \hwunread(1.75,15.5){○}
  \hw(1.0,16.3){2\ 1}

  % 21|21 と「>」、コピーをまたぐ弧
  \hw(3.5,15.3){21}
  \hwline(4.35,14.9)(4.35,15.9)
  \hw(4.45,15.3){21}
  \hwline(3.75,16.0)(4.0,15.95) \hwline(4.0,15.95)(3.85,16.25)
  \hwline(4.95,15.95)(5.2,15.9) \hwline(5.2,15.9)(5.05,16.2)
  \draw[ink,hand,line width=0.7pt,-{Stealth[length=1.8mm]}] (3.55,15.75) to[bend right=40] (5.3,16.5);
  \redpen(0.6,13.9){\small サンプル 1: コピー内 2 + コピー間 1 = 3}

  % 上のジグザグ（B = 2,1,2,1）
  \hw(6.2,14.45){\normalsize 2}
  \hw(7.6,15.75){\normalsize 1}
  \hw(7.95,13.85){\normalsize 2}
  \hw(8.8,15.85){\normalsize 1}
  \draw[ink,hand,line width=0.7pt] (6.85,14.4) -- (7.95,15.25) -- (8.25,14.35) -- (8.9,15.2);
  \foreach \p in {(6.85,14.4),(7.95,15.25),(8.25,14.35)} {\draw[ink,line width=0.6pt] \p circle (0.06);}
  \fill[ink] (8.9,15.2) circle (0.05);
  \draw[ink,hand,line width=0.7pt,-{Stealth[length=1.8mm]}] (6.85,14.6) -- (8.8,14.95);
  \draw[ink,hand,line width=0.7pt,-{Stealth[length=1.6mm]}] (7.05,14.1) to[bend left=35] (7.55,14.75);
  \draw[ink,hand,line width=0.7pt] (8.4,14.1) to[bend left=35] (8.95,14.6);

  % サンプル 2
  \hw(1.4,17.6){3\ 5}
  \hw(1.5,18.3){1\ \ 1\ \ 1}
  \hw(4.0,17.35){1 -- 1}
  \hw(6.5,17.15){0}
  \redpen(0.6,19.4){\small 同じ値どうしは転倒にならない（等号は数えない）→ 0}

  % 下のジグザグ（弧つき）
  \draw[ink,hand,line width=0.7pt] (8.2,16.75) -- (9.05,17.4) -- (9.8,16.75) -- (10.9,17.25) -- (11.6,16.5) -- (12.6,17.15);
  \foreach \p in {(9.8,16.75),(10.9,17.25),(11.6,16.5)} {\draw[ink,line width=0.6pt] \p circle (0.06);}
  \draw[ink,line width=0.6pt] (9.05,17.4) circle (0.06);
  \draw[ink,hand,line width=0.6pt] (9.05,17.4) circle (0.22);
  \draw[ink,hand,line width=0.7pt,-{Stealth[length=1.8mm]}] (8.3,16.8) -- (12.0,17.05);
  \draw[ink,hand,line width=0.7pt] (8.35,16.9) -- (11.85,17.15);
  \draw[ink,hand,line width=0.7pt] (8.8,16.7) to[bend left=35] (9.05,17.1);
  \draw[ink,hand,line width=0.7pt] (9.9,16.65) to[bend left=35] (10.2,17.2);
  \draw[ink,hand,line width=0.7pt] (10.9,16.9) to[bend left=35] (11.3,17.0);
  \draw[ink,hand,line width=0.7pt] (12.2,16.8) to[bend left=35] (12.6,16.95);
  \draw[ink,hand,line width=0.6pt] (9.25,17.2) to[bend left=25] (9.4,17.95);
  \hwunread(13.1,16.5){点の値}
  \redpen(9.6,18.6){\small 合計の式はメモに無く、コードで直接書いた\\\small $x\frac{K(K+1)}{2} + y\frac{K(K-1)}{2}$ → 1 回目で AC}
\end{memopage}

% ---------------------------------------------------------------- 振り返り
\clearpage
\section*{jsc2019\_qual B --- Kleene Inversion}

\begin{itemize}
  \item 問題: \url{https://atcoder.jp/contests/jsc2019-qual/tasks/jsc2019_qual_b}
  \item 提出（2026-08-04, Python）:
    \href{https://atcoder.jp/contests/jsc2019-qual/submissions/78120054}{AC}（1 回目で AC）。コードは \texttt{other/jsc/2019/b/main.py}
  \item 手書き原本: \texttt{Box/Photo Backup/手書き/atcoder/jsc2019\_qual/b/}（\texttt{20260925\_174814.jpg} の 1 枚。撮影 2026-09-25）。
    前の 1 ページが再現。赤は後から入れた振り返り
\end{itemize}

\subsection*{メモで考えたこと}
\begin{enumerate}
  \item 長さ $N$ の $A$ を $K$ 回つなげた $B$（長さ $KN$）の転倒数 $t$ を $10^9+7$ で割った余り、と書き直した。
    例: $A = 1,3,2$, $K=2$ なら $B = 1,3,2,1,3,2$。
  \item $t = \mathrm{count}\{ (i,j) : 0 \leq i < j \leq KN-1,\ B_i > B_j \}$ と式にした。
  \item サンプル \verb|2 2 / 2 1|（$B = 2,1,2,1$）で、大きい値から小さい値へ弧を描いて転倒ペアを数えた。
    同じコピーの中のペアと、コピーをまたぐペアを描き分けている。
\end{enumerate}

\subsection*{つまずき}
無し。弧の図で「同じコピー内」と「コピー間」を分けて見たのが、そのまま解法になった。

\subsection*{正しい考え方}
$A$ の中で $i<j$ かつ $A_i > A_j$ の数を $x$、$i<j$ かつ $A_i < A_j$ の数を $y$ とする。

\begin{itemize}
  \item 同じコピー内の転倒: $K$ 個のコピーそれぞれに $x$ 個 → $Kx$。
  \item コピー間（前のコピーの $A_i$ と後ろのコピーの $A_j$）: 位置に関係なく $A_i > A_j$ の組が転倒になる。
    そういう $(i,j)$ の数は $x + y$、コピーの組は $K(K-1)/2$。
\end{itemize}

合計
\[
  Kx + (x+y)\frac{K(K-1)}{2} = x\frac{K(K+1)}{2} + y\frac{K(K-1)}{2}
\]
提出コードはこの形（サンプル 3 例とも一致: 3, 0, 185297239）。

サンプル 1（$B = 2,1 \mid 2,1$）で内と間を分けると次のとおり。黒はコピー内、赤はコピーをまたぐ転倒。

\begin{center}
\begin{tikzpicture}[x=1.1cm,y=1cm]
  % B = 2,1 | 2,1
  \foreach \i/\v in {0/2,1/1,2/2,3/1} {\node (b\i) at (\i,0) {\large $\v$};}
  \draw[dashed,gray] (1.5,-0.5) -- (1.5,1.6);
  \node[gray,below] at (0.5,-0.4) {\small コピー 1};
  \node[gray,below] at (2.5,-0.4) {\small コピー 2};
  \draw[-{Stealth},thick] (b0.north) to[bend left=60] node[above,font=\small] {内} (b1.north);
  \draw[-{Stealth},thick] (b2.north) to[bend left=60] node[above,font=\small] {内} (b3.north);
  \draw[-{Stealth},thick,redpen] (b0.north) to[bend left=50] node[above,font=\small] {間} (b3.north);
  \node[anchor=west,font=\small] at (4.2,1.0) {内: $Kx = 2 \cdot 1 = 2$};
  \node[anchor=west,font=\small] at (4.2,0.4) {間: $(x+y)\frac{K(K-1)}{2} = 1 \cdot 1 = 1$};
  \node[anchor=west,font=\small] at (4.2,-0.2) {合計 3};
\end{tikzpicture}
\end{center}

\subsection*{次に活かすこと}
\begin{itemize}
  \item 繰り返し構造の数え上げは、\textbf{ブロック内} と \textbf{ブロック間} に分けると、ブロック間は位置の制約が消えて数えやすくなる。
\end{itemize}

\end{document}
